%Vision-language models can describe scientific figures fluently, yet often behave unreliably when evidence is missing or misleading. Existing benchmarks focus on perception and reasoning accuracy while overlooking behavioural reliability under uncertainty. We introduce \textsc{SciFig-Eval}, a diagnostic benchmark for scientific figure understanding that evaluates perception, reasoning, and behavioural robustness through controlled probes.
%Vision-language models can generate fluent descriptions of scientific figures, yet often behave unreliably when visual evidence is missing, ambiguous, or misleading. Existing benchmarks primarily measure perception and reasoning accuracy, while overlooking behavioural reliability under uncertainty. We present \textsc{SciFig-Eval}, a diagnostic benchmark for scientific figure understanding that evaluates perception, reasoning, and robustness through controlled behavioural probes. 
%The benchmark contains 250 annotated figures and more than 23,000 evaluations across 8 models, combining MQM-based description scoring, 1000 targeted capability questions, image transformations, misleading captions, false-premise probes, and selective-blur tests. We further introduce the Admittance-Resistance-Inductance (A-R-I) framework to measure whether models acknowledge insufficient evidence, resist misleading context, and make cautious inferences from partial visual information. Our experiments reveal substantial behavioural differences among models with similar perception performance. GPT-5.2 achieves the highest description quality (MQM 91.6) but hallucinates answers about unreadable elements in 98\% of cases, whereas Gemini 3.1 Pro attains comparable perception performance (MQM 90.2), admits uncertainty in 90\% of such cases, and achieves the strongest resistance score (0.91). These results show that high perception accuracy alone does not guarantee reliable scientific reasoning, highlighting the need for evaluation frameworks that explicitly measure uncertainty handling and robustness.
% Vision-language models can describe scientific figures fluently, yet often behave unreliably when evidence is missing or misleading. 
Existing vision-language model (VLM) benchmarks emphasize perception and reasoning accuracy (\textit{how well VLMs describe and reason about what they see in an image}), with limited attention to behavioral reliability under uncertainty (\textit{how they behave when visual evidence is missing or misleading}). We introduce \textsc{SciFig-Eval}, a diagnostic VLM benchmark for scientific figure understanding that jointly evaluates perception, reasoning, and behavioral reliability under uncertainty. It contains 250 figures with high-quality human annotations in three evaluation aspects outlined, totaling 600+ hours of annotation effort. We further extend these figures via image transformations,
reasoning questions, 
resistance probes, 
caption-bias probes, and 
confirmed selective-blur targets, producing over 34,000 evaluation setups for stress testing.
We further propose the Admittance--Resistance--Inductance (A-R-I) framework to evaluate whether models acknowledge insufficient evidence, resist misleading context, and infer cautiously from partial information. Our results reveal substantial behavioral differences among models. GPT-5.2 achieves the highest description quality (MQM 91.6) with strong reasoning accuracy (78.4\%), yet hallucinates unreadable content in 96\% of cases, whereas Gemini 3.1 Pro, a comparably capable model (MQM 90.2, reasoning 81.0\%), admits uncertainty in 71\% of such cases and achieves the strongest resistance score (0.91). These findings show that \textit{high perception and reasoning accuracy alone do not guarantee behavioral reliability, a dimension critical for deployment in scientific workflows.
}
% a dimension critical for VLMs}. %Code and data will be released in the final version.

%Vision-language models can describe scientific figures well while behaving unreliably when evidence is missing or misleading. We present \textsc{SciFig-Eval}, a diagnostic benchmark for scientific figure understanding that evaluates perception, reasoning, and behaviour rather than collapsing them into a single accuracy score. The benchmark contains 250 annotated arXiv figures and more than 23,000 model-output evaluations across eight models, combining checklist-based MQM scoring, 1,000 targeted capability questions, image transformations, modified captions, false-premise probes, and selective-blur tests. We introduce the Admittance--Resistance--Inductance (A-R-I) framework to measure whether models acknowledge unreadable evidence, resist misleading context, and infer only when partial visual information supports it. The results reveal sharp behavioural divergences among models with similar perception scores. GPT-5.2 achieves the best description quality (MQM 91.6) but fabricates answers about unreadable elements in 98\% of cases. Gemini 3.1 Pro scores similarly on perception (90.2), yet admits uncertainty in 90\% of such cases and leads on resistance (0.91). These findings show that perception and reasoning benchmarks alone can miss deployment-critical failure modes in scientific workflows.
